\section{Reduction to the middle degree}\label{sec:reduction}

The two preceding sections dealt with restriction to hypersurfaces and with
the one classical correspondence that avoids it. Here we combine the first of
them with the Lefschetz decomposition to locate the difficulty of the Hodge
conjecture in general. The result, \Cref{thm:equivalence}, is an equivalence
in which neither side is easier than the other. Its use is to isolate the
primitive classes of middle degree as the one place where a new method is
needed, and by \Cref{thm:hyperplane} restriction to hyperplane sections sees
nothing there.

\subsection{Reduction to primitive classes}

\begin{proposition}[Lefschetz decomposition of Hodge classes]%
\label{prop:lefschetzdecomp}
Let $X$ be smooth projective of dimension $d$ with ample class $L$. Every
$\alpha\in\Hdg^{p}(X)$ can be written
\[
  \alpha=\sum_{r\ge\max(0,\,2p-d)}L^{r}\beta_{r},
  \qquad
  \beta_{r}\in\Hdg^{p-r}(X)\cap H^{2p-2r}_{\mathrm{prim}}(X,\QQ),
\]
and the $\beta_{r}$ are uniquely determined. Consequently the Hodge conjecture
holds for $X$ if and only if every primitive rational Hodge class on $X$ is
algebraic.
\end{proposition}

\begin{proof}
The Lefschetz decomposition
$H^{k}(X,\QQ)=\bigoplus_{r}L^{r}H^{k-2r}_{\mathrm{prim}}(X,\QQ)$ is a
decomposition of rational Hodge structures, because $L$ is a rational class
of type $(1,1)$ and the primitive parts are kernels of morphisms of Hodge
structures. The rationality of $L$ is essential here: for a K\"ahler class
that is not a real multiple of a rational one the primitive subspaces are in
general not rational, and a rational Hodge class need not decompose into
rational primitive pieces. A Hodge class $\alpha$ therefore has Hodge classes
$L^{r}\beta_{r}$ as its components. By hard Lefschetz, $L^{r}$ is injective
on $H^{2p-2r}_{\mathrm{prim}}(X,\QQ)$ for $r\ge\max(0,2p-d)$, and it is a
morphism of Hodge structures of bidegree $(r,r)$, so each $\beta_{r}$ is
uniquely determined and is itself a Hodge class. For the last statement, if
each primitive $\beta_{r}$ is algebraic then so is $L^{r}\beta_{r}$, since $L$
is algebraic and algebraic classes form a subring, and hence so is $\alpha$;
the converse holds because primitive Hodge classes are Hodge classes.
\end{proof}

\subsection{The equivalence}

\begin{theorem}\label{thm:equivalence}
The following are equivalent.
\begin{enumerate}[label=\textup{(\roman*)}]
\item The Hodge conjecture holds for every smooth complex projective variety.
\item The Lefschetz standard conjecture $B(X)$ holds for every smooth complex
  projective variety $X$, and every primitive rational Hodge class of middle
  degree on every smooth complex projective variety of even dimension is
  algebraic.
\end{enumerate}
\end{theorem}

\begin{proof}
(i) implies (ii). The second half of (ii) is a special case of (i). For
$B(X)$, note that $\Lambda$ is a morphism of Hodge structures of bidegree
$(-1,-1)$, hence is given by a Hodge class on $X\times X$ through the
standard correspondence between morphisms of Hodge structures and Hodge
classes on the product, and (i) applied to $X\times X$ makes that class
algebraic. This is Kleiman's argument \cite[\S2]{Kle68}.

(ii) implies (i). We argue by induction on $d=\dim X$. For $d\le1$ the Hodge
conjecture is classical. Let $d\ge2$ and assume the conjecture for all smooth
projective varieties of dimension less than $d$. By Bertini there is a smooth
hypersurface $j\colon Y\hookrightarrow X$ in $|mH|$ for an ample $H$ and
$m\gg0$; put $L=c_{1}(\cO_{X}(Y))$, and take primitivity and the operator
$\Lambda$ with respect to this $L$, which is harmless because $B(X)$ does not
depend on the choice of ample class \cite{Kle68,Kle94}. By
\Cref{prop:lefschetzdecomp} it suffices to prove that a primitive Hodge class
$\alpha\in H^{2p}(X,\QQ)$ is algebraic and, by hard Lefschetz, primitive
classes occur only in degrees at most $d$, so $2p\le d$.

If $2p=d$, then $\alpha$ is a primitive Hodge class of middle degree on a
variety of even dimension, so it is algebraic by the second half of (ii).

If $2p<d$, the class $j^{*}\alpha$ is a rational Hodge class on $Y$ because
$j^{*}$ is a morphism of Hodge structures. As $\dim Y=d-1<d$, the inductive
hypothesis makes $j^{*}\alpha$ algebraic, say $j^{*}\alpha=\cl(W)$ with
$W\in\CH^{p}(Y)_{\QQ}$. Then
$\cl(j_{*}W)=j_{*}j^{*}\alpha=L\alpha$ by \Cref{lem:gysinrestrict}, so
$L\alpha$ is algebraic. By $B(X)$ the operator $\Lambda$ carries algebraic
classes to algebraic classes, and by \Cref{lem:LambdaL},
$\alpha=(d-2p)^{-1}\Lambda(L\alpha)$, which is therefore algebraic; the scalar
is defined because $d-2p\ge1$.
\end{proof}

\begin{remark}[The content of the equivalence]\label{rem:notsimpler}
$B(X)$ is itself a consequence of (i), so
\Cref{thm:equivalence} does not reduce the problem to a smaller one; it
locates the difficulty. The induction in the proof that (ii) implies (i) runs
through every degree except the middle one, where it rests entirely on the
hypothesis. By \Cref{thm:hyperplane}, the step it uses below the middle
degree, restriction to a hypersurface, returns zero on primitive classes of
middle degree, so no refinement of that step extends the induction to that
degree.
\end{remark}

\begin{corollary}[Primitive classes of middle degree]\label{cor:whereidea}
Any proof of the Hodge conjecture must, at some point, produce an algebraic
cycle representing a primitive Hodge class of middle degree on a variety of
even dimension, by a method that does not factor through restriction to
hyperplane sections of that variety.
\end{corollary}

\begin{proof}
By \Cref{thm:equivalence}, the Hodge conjecture includes the algebraicity of
the primitive Hodge classes of middle degree on varieties of even dimension,
and by \Cref{cor:blind} every such class restricts to zero on every smooth
ample hypersurface.
\end{proof}

\begin{remark}[Methods that do not factor through hyperplane sections]%
\label{rem:othermethods}
\Cref{cor:whereidea} does not say that the conjecture is out of reach, and
the settled cases come from known methods outside its scope. The
Kuga--Satake construction (\Cref{sec:ks}) replaces a weight-two Hodge
structure of K3 type by the $H^{1}$ of an abelian variety; the class it
produces is of type $(1,1)$ on a product, hence algebraic by Lefschetz, and it
is not a restriction to a subvariety. Degenerations and limits of Hodge
structures move the variety instead of cutting it. Correspondences from other
varieties, of which the Kuga--Satake construction is one instance and
Markman's construction \cite{Mar25a,Mar25b} another, are also outside the
scope. Arguments through the coniveau filtration, which show that a class
supported on a proper closed subvariety is controlled by that subvariety, are
compatible with everything here: they lower the dimension legitimately,
because the class is supported on the smaller variety, whereas a primitive
class of coniveau zero is supported on no proper closed subvariety.
\end{remark}

